\documentclass[12pt]{article}
\usepackage[margin=1in]{geometry}
\usepackage{amsmath}
\usepackage{parskip}
\usepackage{hyperref}
\usepackage{enumitem}
\usepackage{titlesec}

\title{Advanced NLP -- Exercise 1 \\ \large Open Questions}
\author{Netanel Azran}
\date{}

\begin{document}
\maketitle

\noindent\textbf{GitHub Repository:} \url{https://github.com/netanelazran11/ANLP-2026}

\section*{Question 1: QA Datasets Measuring Intrinsic Language Properties}

Below are three QA datasets that use the question-answering format to annotate intrinsic
properties of language understanding -- that is, properties of language itself rather than
performance on some downstream application.

\begin{enumerate}

    \item \textbf{bAbI} (Weston et al., 2015)

    bAbI consists of 20 synthetic QA tasks, each isolating a specific reasoning ability
    such as spatial reasoning, coreference, or multi-hop fact chaining. Since every task
    is designed to test one elementary language faculty in isolation (e.g., tracking entity
    locations across sentences), it directly measures intrinsic comprehension skills rather
    than performance on any downstream application.

    \item \textbf{Winograd Schema Challenge (WSC)} (Levesque, 2012)

    WSC uses a QA format to test pronoun coreference resolution: given a sentence with an
    ambiguous pronoun, the model must identify the correct referent using world knowledge
    rather than surface-level cues. This is an intrinsic task because resolving anaphora
    is a core linguistic competence -- the dataset is explicitly designed to measure this
    ability, not to solve any real-world application.

    \item \textbf{DROP} (Dua et al., 2019)

    DROP is a reading comprehension dataset where answering requires discrete reasoning
    operations over the passage -- such as counting, arithmetic, or sorting -- using a
    QA format. It measures an intrinsic property of language understanding because the
    questions are designed to probe symbolic reasoning over text rather than surface-level
    pattern matching, directly testing whether a model can compose and apply reasoning
    steps from linguistic input.

\end{enumerate}

\section*{Question 2: Inference-Time Scaling Methods}

\subsection*{(a) Methods Overview}

\begin{enumerate}

    \item \textbf{Self-Consistency} (Wang et al., 2022)

    \textbf{Description:} The model samples $N$ independent Chain-of-Thought reasoning
    paths for the same prompt, and the final answer is chosen by majority vote.

    \textbf{Advantages:} Improves over single greedy decoding without requiring any
    additional model or training; naturally reduces variance across reasoning paths.

    \textbf{Computational bottleneck:} Requires generating $N$ full CoT chains, so
    compute and memory scale with both $N$ and chain length.

    \textbf{Parallelizable:} Yes -- all $N$ chains are independent and can be batched
    in a single forward pass.

    \item \textbf{Verifiers at Test Time} (Cobbe et al., 2021; Lightman et al., 2023)

    \textbf{Description:} A separate verifier model scores multiple candidate outputs
    generated by the model, and the highest-scoring one is returned as the final answer.

    \textbf{Advantages:} Can recover correct answers even when the generator alone fails;
    process reward models (PRMs) add fine-grained step-level scoring for better selection.

    \textbf{Computational bottleneck:} Requires both $N$ generation passes and $N$
    verifier passes; PRMs are especially expensive as they score each reasoning step.

    \textbf{Parallelizable:} Yes -- generation and scoring are independent across
    candidates and can be fully batched.

    \item \textbf{CoT Length Scaling / Thinking Models} (OpenAI o1, 2024; DeepSeek R1, 2025)

    \textbf{Description:} The model is trained to produce extended reasoning chains that
    include planning, backtracking, and self-evaluation before emitting a final answer.

    \textbf{Advantages:} Allows the model to correct mistakes mid-generation and recover
    from dead ends; achieves strong results on hard scientific and mathematical tasks.

    \textbf{Computational bottleneck:} Token generation is strictly sequential, so longer
    chains directly increase latency and KV cache memory usage.

    \textbf{Parallelizable:} No -- each token is conditioned on all previous tokens,
    making generation within a single chain inherently sequential.

\end{enumerate}

\subsection*{(b) Method Choice for a Complex Scientific Reasoning Task (Single GPU, Large Memory)}

I would choose \textbf{Self-Consistency}.

Since the task requires reasoning, a single greedy decoding pass is unreliable -- the
model may follow an incorrect reasoning path and commit to a wrong answer. Self-Consistency
addresses this by sampling many independent CoT chains and taking a majority vote, which
consistently outperforms single-pass decoding on reasoning benchmarks.

Given a single GPU with large memory, all $N$ chains can be batched into a single forward
pass, directly exploiting the available memory. CoT Length Scaling is inherently sequential
and cannot benefit from batching, while Verifiers require a second model that competes for
the same GPU memory. Self-Consistency is therefore the most effective choice under these
hardware constraints.

\section*{Qualitative Analysis: Best vs. Worst Configuration}

\subsection*{Did the best validation configuration also achieve the best test accuracy?}

No. After running all three trained checkpoints on the test set, the results were:

\begin{center}
\begin{tabular}{lcc}
\hline
\textbf{Configuration} & \textbf{Val Acc} & \textbf{Test Acc} \\
\hline
ep=3, lr=2e-5, bs=16 & 85.54\% & 82.96\% \\
ep=4, lr=3e-5, bs=32 & 84.80\% & \textbf{83.25\%} \\
ep=2, lr=1e-5, bs=16 & 78.68\% & 76.41\% \\
\hline
\end{tabular}
\end{center}

The configuration with the highest validation accuracy (ep=3, lr=2e-5, bs=16) did
\textbf{not} achieve the highest test accuracy. The second configuration (ep=4, lr=3e-5,
bs=32) scored slightly higher on the test set (83.25\% vs.\ 82.96\%), despite lower
validation accuracy. This suggests mild overfitting to the validation set by the
best-validation model, and shows that validation accuracy is not always a perfect proxy
for generalization.

\subsection*{Error Analysis}

We compared the best configuration (3 epochs, lr=2e-5, bs=16, val acc=85.5\%) against
the worst (2 epochs, lr=1e-5, bs=16, val acc=78.7\%) by running both models on the
validation set and identifying examples where the best model predicted correctly but the
worst model failed.

Nearly all such cases shared the same pattern: the true label was \textbf{0 (not a
paraphrase)}, but the worst model incorrectly predicted \textbf{1 (paraphrase)}. In other
words, the worst model was systematically fooled by sentence pairs that \textit{look}
similar on the surface but differ in meaning or completeness.

A few representative examples illustrate this:

\begin{itemize}

    \item \textit{``No dates have been set for the civil or the criminal trial.''} vs.
    \textit{``No dates have been set for the criminal or civil cases, but Shanley has
    pleaded not guilty.''} — the sentences share most of their words, but S2 adds a
    semantically distinct clause that makes them non-paraphrases.

    \item \textit{``Sanitation is poor\ldots there could be typhoid and cholera''} vs.
    the same sentence with an extra clause about drinking water inserted — again, high
    lexical overlap but different content.

    \item \textit{``Danbury prosecutor Warren Murray could not be reached for comment
    Monday.''} vs. \textit{``Prosecutors could not be reached for comment after the legal
    papers were obtained late Monday afternoon.''} — same general event, but different
    specificity (named individual vs. generic ``prosecutors'').

\end{itemize}

The worst model, trained for only 2 epochs with a very low learning rate, did not converge
sufficiently to distinguish semantic equivalence from surface-level lexical overlap. It
defaulted to predicting paraphrase whenever sentences shared a large fraction of their
vocabulary. The best model, trained for 3 epochs with a higher learning rate, learned to
attend to the specific details that differentiate near-miss pairs -- such as added clauses,
different levels of specificity, or subtle changes in who is doing what -- and therefore
correctly classified them as non-paraphrases.

\end{document}
